\section{System Architecture and Methodology}
\label{sec:methodology}

Our system is organised as a four-phase pipeline: multi-source data
collection, feature engineering, hybrid neural-network training, and
rigorous evaluation.  Figure~\ref{fig:architecture} provides an overview
of the complete architecture.

% ─────────────────────────────────────────────────────────────────────────
\subsection{Phase 1: Multi-Source Data Collection}
\label{subsec:collection}

\subsubsection{Seed Arrays and Transaction Neighbourhood}
The pipeline starts from two explicit Python arrays stored in
\texttt{code/data/addresses/}:

\begin{itemize}
    \item \texttt{BENIGN\_ADDRESSES}: 504 verified blue-chip Ethereum
    contracts (ERC-20 tokens, DEX pools, lending protocols) sourced from
    the CoinGecko top-N list and cross-validated against the Etherscan
    verified-contract registry.

    \item \texttt{FRAUD\_ADDRESSES}: 156 victim-protocol contracts
    extracted from DeFiHackLabs proof-of-concept (PoC) Solidity
    files~\cite{defihacklabs2024}.  Each address is identified by a
    \texttt{// Vulnerable Contract : https://etherscan.io/} comment in
    the PoC, verified to carry Solidity source code on Etherscan, and
    confirmed to have on-chain transaction history---properties that are
    essential for training both the graph encoder and the language
    encoder.
\end{itemize}

From the union of these two arrays we retrieve the transaction
neighbourhood from Google BigQuery's public Ethereum
dataset~\cite{bigquery2024}.  Our balanced sampling strategy allocates
up to 2,000 most-recent transactions per fraud address and up to 300 per
benign address, yielding a dense, label-rich graph of \num{239742}
transactions across \num{85680} unique addresses while avoiding the
extreme volume dominance of high-traffic benign contracts such as WETH
and USDT.

\subsubsection{Ground-Truth Labels}
\label{subsubsec:labels}

Our fraud labels come from two sources with deliberately different roles.

\textbf{Primary source: DeFiHackLabs victim contracts.}
Of the 156 DeFiHackLabs seed addresses, 130 appear in the transaction
neighbourhood (26 were not captured by the balanced sampling); these
contribute \textbf{121 of the 128 fraud-labelled nodes} (94.5\%) in the
final graph.  Because each address was selected from a PoC repository that
documents real exploits with on-chain evidence, these labels carry a very
high degree of confidence.

\textbf{Secondary augmentation: Forta Network labels.}
We cross-reference all graph nodes against the Forta Network labelled-datasets
release~\cite{fortalabels}, which provides community-verified fraud
annotations covering phishing, Ponzi schemes, and malicious smart
contracts.  This adds \textbf{7 additional fraud-labelled nodes} (5.5\%)
not already covered by the DeFiHackLabs seeds.  All fraud sub-categories
are collapsed into a single label ($y=1$).

All remaining addresses receive $y=0$ (benign), subject to the
closed-world assumption discussed in Section~\ref{subsec:limitations}.

Table~\ref{tab:dataset} summarises the complete dataset statistics.

\begin{table}[t]
\centering
\caption{Dataset statistics after full pipeline execution.}
\label{tab:dataset}
\begin{tabular}{lr}
\toprule
Statistic & Value \\
\midrule
\multicolumn{2}{l}{\textit{Seed arrays}} \\
\ \ Fraud seeds (DeFiHackLabs victim contracts) & 156 \\
\ \ Benign seeds (blue-chip protocol contracts) & 504 \\
\ \ Total seed addresses & 660 \\
\midrule
\multicolumn{2}{l}{\textit{Transaction graph}} \\
\ \ Total nodes (unique addresses) & 85,680 \\
\ \ Total edges (transactions) & 239,742 \\
\midrule
\multicolumn{2}{l}{\textit{Fraud labels}} \\
\ \ From DeFiHackLabs seeds & 121 (94.5\%) \\
\ \ From Forta augmentation only & 7 (5.5\%) \\
\ \ Total fraud-labelled nodes & 128 \\
\midrule
\multicolumn{2}{l}{\textit{Eligible for training (label + source code)}} \\
\ \ Fraud contracts & 95 \\
\ \ Benign contracts & 209 \\
\ \ Total eligible & 304 \\
\ \ Training / Validation / Test & 212 / 45 / 47 \\
\bottomrule
\end{tabular}
\end{table}

\subsubsection{Source Code Retrieval}
For each unique address in the neighbourhood, we query the Etherscan
API v2~\cite{etherscanapi} for verified Solidity source code.  Source
is successfully retrieved for \num{304} addresses; the remainder are
externally-owned accounts or unverified contracts.  All source strings
are Base64-encoded for safe CSV storage and later decoded at tokenisation
time.

% ─────────────────────────────────────────────────────────────────────────
\subsection{Phase 2: Feature Engineering}
\label{subsec:features}

\subsubsection{Graph Integration via DuckDB}
All three data frames (transactions, Forta labels, Etherscan source) are
registered as views in an in-process DuckDB~\cite{duckdb2019} instance.
A sequence of \texttt{LEFT JOIN} operations materialises two output
tables: \texttt{nodes} (address, contract\_text, label $y$, is\_labeled)
and \texttt{edges} (src, dst, value, tx\_hash).  All addresses are
lower-cased before any join, eliminating the silent mismatches that arise
from EIP-55 checksummed addresses in Forta exports versus lowercase
BigQuery output.

\subsubsection{Node Feature Engineering}
\label{subsubsec:features}

For each node $v$ we compute eight features that capture both structural
position and transactional value flow:

\begin{equation}
\mathbf{f}_v = \bigl[
  d^{\text{in}}_v,\;
  d^{\text{out}}_v,\;
  d^{\text{in}}_v + d^{\text{out}}_v,\;
  \bar{w}^{\text{in}}_v,\;
  \bar{w}^{\text{out}}_v,\;
  \log(1 + W_v),\;
  u^{\text{in}}_v,\;
  u^{\text{out}}_v
\bigr]^\top \in \mathbb{R}^8,
\label{eq:features}
\end{equation}

where $d^{\text{in}}_v$ and $d^{\text{out}}_v$ are in- and out-degrees;
$\bar{w}^{\text{in}}_v$ and $\bar{w}^{\text{out}}_v$ are the mean
incoming and outgoing transaction values (in Wei); $W_v$ is the total
Wei volume; and $u^{\text{in}}_v$, $u^{\text{out}}_v$ are the counts of
unique counterparty addresses for incoming and outgoing transactions,
respectively.  Value features are log-compressed to reduce the
multi-order-of-magnitude spread in Ethereum transaction sizes.

These features encode complementary fraud signals: phishing addresses
typically exhibit high $d^{\text{in}}_v$ (many victims) relative to
$d^{\text{out}}_v$ (few recipients), whereas Ponzi contracts show
characteristic value-flow patterns~\cite{liang2025ponziguard}.

\subsubsection{Source-Code Preprocessing}
Solidity source strings are decoded from Base64 and tokenised using the
CodeBERT tokenizer with a maximum length of 512 tokens.  Sequences
exceeding the limit are truncated; shorter sequences are padded.  This
pre-processing is performed lazily inside \texttt{ContractDataset}
at batch-construction time.

% ─────────────────────────────────────────────────────────────────────────
\subsection{Phase 3: Hybrid Neural Network Architecture}
\label{subsec:network}

Figure~\ref{fig:fusion} illustrates the internal architecture of the
fusion classifier.

\subsubsection{Graph Encoder}
A two-layer GraphSAGE network~\cite{hamilton2017graphsage} maps the
8-dimensional node features into a 64-dimensional latent space:

\begin{align}
\mathbf{h}_v^{(1)} &= \sigma\!\Bigl(
  \mathbf{W}_1 \cdot \mathrm{MEAN}\!\bigl(\{\mathbf{f}_u : u \in \mathcal{N}(v)\cup\{v\}\}\bigr)
\Bigr), \label{eq:sage1} \\
\mathbf{h}_v^{(2)} &= \mathbf{W}_2 \cdot \mathrm{MEAN}\!\bigl(
  \{\mathbf{h}_u^{(1)} : u \in \mathcal{N}(v)\cup\{v\}\}\bigr),
  \label{eq:sage2}
\end{align}

where $\sigma$ is ReLU, $\mathbf{W}_1 \in \mathbb{R}^{128 \times 8}$,
and $\mathbf{W}_2 \in \mathbb{R}^{64 \times 128}$.  Dropout($p=0.5$) is
applied after the first layer.  GraphSAGE is \emph{inductive}: it learns
aggregation functions that generalise to nodes unseen during training,
which is essential for screening newly deployed contracts.

\subsubsection{Language Encoder}
We load \texttt{microsoft/codebert-base}~\cite{feng2020codebert} from
Hugging Face Transformers, freeze all 125M parameters, and extract the
768-dimensional \texttt{pooler\_output} as the contract representation.
Freezing CodeBERT bounds training memory within a single GPU and isolates
the fusion head as the sole trainable bridge between modalities, allowing
the model to learn how to combine---rather than re-learn---the pre-trained
code semantics.  Prior work on fine-tuning LLMs for classification
tasks~\cite{poudel2026fakenews} demonstrates that frozen pre-trained
representations already provide strong signals, validating this design
choice; parameter-efficient fine-tuning (e.g.\ LoRA~\cite{hu2022lora})
is deferred to future work.

\subsubsection{Fusion Classifier}
\label{subsubsec:fusion}

The 64-dimensional graph embedding $\mathbf{h}_v^{\text{GNN}}$ and the
768-dimensional code embedding $\mathbf{h}_v^{\text{LLM}}$ are
concatenated to form an 832-dimensional joint representation.  The
classification head then computes:

\begin{equation}
\hat{y}_v = \mathbf{W}_4 \, \sigma\!\Bigl(
  \mathbf{W}_3 \,[\mathbf{h}_v^{\text{GNN}};\, \mathbf{h}_v^{\text{LLM}}] + \mathbf{b}_3
\Bigr) + \mathbf{b}_4,
\label{eq:fusion}
\end{equation}

where $\mathbf{W}_3 \in \mathbb{R}^{128 \times 832}$ is the fusion
projection, $\mathbf{W}_4 \in \mathbb{R}^{2 \times 128}$ is the
classifier, and Dropout($p=0.5$) is inserted before each linear layer.
Total trainable parameter count is \num{124226}; the remaining
$\sim$124.6M parameters belong to frozen CodeBERT.

The concatenation bottleneck forces the model to discover a compact joint
embedding rather than memorising individual modalities.
Section~\ref{subsec:ablation_module} compares this with an
attention-based fusion alternative.

% ─────────────────────────────────────────────────────────────────────────
\subsection{Phase 4: Training Strategy}
\label{subsec:training}

\subsubsection{Class-Weighted Loss}
Of the 304 labelled contracts, 95 are fraud (31.3\%).  We address the
resulting imbalance with weighted cross-entropy:
\begin{equation}
\mathcal{L} = -\sum_v \bigl[
  w_0(1-y_v)\log\hat{p}_{v,0} + w_1 y_v \log\hat{p}_{v,1}
\bigr],
\label{eq:loss}
\end{equation}
where $w_0 = 1.0$ and $w_1 = n_{\text{neg}} / n_{\text{pos}} \approx
2.37$ (computed from the training split).

\subsubsection{Optimiser and Scheduler}
We use AdamW~\cite{loshchilov2019adamw} with learning rate $10^{-3}$ and
weight decay $10^{-4}$.  A cosine annealing
schedule~\cite{loshchilov2017cosine} decays the learning rate to
$10^{-5}$ over a 100-epoch budget.  Gradient clipping at norm 1.0
suppresses early-epoch instability.

\subsubsection{Early Stopping and Model Selection}
Training halts when the validation fraud F1-score fails to improve for
20 consecutive epochs.  The checkpoint with the highest validation fraud
F1 is reloaded before test-set evaluation; only this checkpoint is
reported.  Our fusion model early-stopped at epoch 41 (best checkpoint at epoch 21) with best
validation F1 = 0.966.

\subsubsection{Data Split}
The 304 eligible contracts (those with both verified source and an
explicit label) are partitioned 70/15/15 into train (212), validation
(45), and test (47) sets using a stratified random split seeded at 42.
The identical split is used for all ablation baselines, ensuring
controlled comparisons.
